\begin{subsection}{A Sheaf on Monetary Risk Measures}
    Let \( \mathcal{X} = L^\infty(\Omega, \Sigma, \mathbb{P}) \) with its usual topology for this section.
    Define a contravariant functor based on the map \( \rejectidb{\cdot} \) from Definition \ref{rejectionsetsarefilters}
    with the signature \( \presh :\textbf{Open}(\mathcal{X}) \to \quotlat{\extrm}, U \mapsto \frac{\extrm}{\rejectidb{U}} \), where we understand 
    \( \textbf{Open}(\mathcal{X}) \) as the category with objects the open sets \( U \subseteq \mathcal{X} \) and 
    and morphisms being inclusions \( U_1 \subseteq U_2 \) of them. The codomain is the category with objects being the quotient lattices of \( \extrm \)
    and morphism induced as in Remark \ref{dualfilterquotient}.
    \begin{theorem} \label{rmpresheaf}
        The contravariant functor \( \presh : \textbf{Open}(\mathcal{X}) \to \quotlat{\extrm} \) defines a separated presheaf on
        \( \mathcal{X} \).
    \end{theorem}
    \begin{proof}
        \begin{itemize}
            \item \( \presh \) is a presheaf.
                \begin{itemize}
                    \item
                        If \( U = \emptyset \) then \( \rejectidb{U} = \extrm \) and \( \frac{\extrm}{\rejectidb{U}} = \{0\} \) the trivial lattice, clearly all lattices have a lattice morphism 
                        to \( \{0\} \).
                        For each open \( U \subseteq \mathcal{X} \) with the identity morphism of sets there's also an identity morphism of lattices \( \frac{\extrm}{\rejectidb{U}} \).
                        For open \( V \subseteq U \) we get by Lemma \ref{rejectionsetsarefilters}, that \( \rejectidb{U} \subseteq  \rejectidb{V} \) as filters in \( \extrm \)
                        and Remark \ref{dualfilterquotient} then gives a morphism \( \varphi_{V,U} : \frac{\extrm}{\rejectidb{U}} \to \frac{\extrm}{\rejectidb{V}}, [\rho]_U \mapsto [\rho]_V \),
                        by Lemma \ref{triangleidealcommutes} it also holds for open \( W \subseteq V \) that \( \varphi_{W,U} = \varphi_{W,V} \circ \varphi_{V,U} \).
                \end{itemize}
            \item \( \presh \) is separated.
                \begin{itemize}
                    \item
                        Take any \( [\rho_1], [\rho_2] \in \presh(U) \), where \( \rho_1, \rho_2 \in \rejectidb{U} \), with \( [\rho_1]|_{U_i} = [\rho_2]|_{U_i} \)
                        for all \( U_i \). This implies for each \( i \) some \( \eta_i \in \rejectidb{U_i} \) s.t. \( \rho_1 \wedge \eta_i = \rho_2 \wedge \eta_i \). 
                        Note that \( \rho_1, \rho_2, \eta_i \in \extrm \), so we can define the risk measure \( \eta = \bigvee_{i} \eta_i \). Since \( \eta \geq \eta_i \) 
                        for all \( i \), it follows \( \eta(X) > 0 \) for all \( X \in U \), so \( \eta \in \rejectidb{U} \). Moreover,
                        \( \rho_1 \wedge \eta = \bigvee_{i} (\rho_1 \wedge \eta_i) = \bigvee_{i} (\rho_2 \wedge \eta_i) = \rho_2 \wedge \eta \), where we used Lemma \ref{extrmisjidandmid}
                        for the equalities, but then \( [\rho_1] = [\rho_2] \).
                \end{itemize}
            
        \end{itemize}
    \end{proof}
    While we were previously working mostly with \( \acceptidealb{U} \), it is quite natural to use \( \rejectidb{U} \) for the presheaf instead.
    First of all acceptance sets are often not open, but closed, while "rejection sets" are typically open. Also we get now a close analogy
    to the local rings. We discovered that every \( X \in \mathcal{X} \) defines an element \( \acceptidealb{\{X\}} \in \text{Spec } \extrm \), but at the same time 
    they also define a filter via \( \rejectidb{\{X\}} \in \text{Spec } \extrm \) and it even holds \( \extrm \setminus \rejectidb{\{X\}} = \acceptidealb{\{X\}} \).
    \begin{proposition}
        \( \acceptidealb{\{X\}} \) is under the canonical map, \( \varphi : \extrm \to \frac{\extrm}{\rejectidb{\{X\}}} \) mapped to the unique maximal ideal
        of \( \frac{\extrm}{\rejectidb{\{X\}}} \).
    \end{proposition}
    \begin{proof}
        By Remark \ref{dualfilterquotient}, \( \varphi(\rejectidb{\{X\}}) = [\top] \), and as \( \acceptidealb{\{X\}} \) is the complement of \( \rejectidb{\{X\}} \)
        all remaining elements of \( \frac{\extrm}{\rejectidb{\{X\}}} \) are in its image.
    \end{proof}
    So even though we did not have maximal ideals for each point \( X \in \mathcal{X} \) in \( \text{Spec } \extrm \), they still define unique maximal ideals in the quotient lattices
    with respect to their filters, as do maximal ideals \( \mathfrak{p} \in \text{Spec } A \) for a commutative ring \( A \) in their local rings \( A_{\mathfrak{p}} \).
    Therefore we could interpret such filters as analogues to the multplicative sets \( A \setminus \mathfrak{p} \) and were granted all that despite not having a concept of inverses in \( \extrm \).
    \newline
    Also we can now answer what the stalks of our presheaf are:
    \begin{corollary}
        The stalk \( \presh_X \) at a point \( X \in \mathcal{X} \) of a presheaf \( \presh \) on the topological space \( \mathcal{X} \),
        is defined as the colimit of the image of the open neighborhoods of the point:
        \[ 
           \presh_{X} := \varinjlim_{x \in U} \presh(U),
        \]
        For the presheaf from Theorem \ref{rmpresheaf} the stalk satisfies
        \[
            \presh_{X} = \frac{\extrm}{\rejectidb{\{X\}}}.
        \]
    \end{corollary}
    \begin{proof}
        By Proposition \ref{colimitforquotientlattice} it suffices to verify 
        that \( \rejectidb{\{X\}} = \cup_{x \in U} \rejectidb{U} \).
        Take \( \rho \in \cup_{x \in U} \rejectidb{U} \), because \( X \in U \), we also have \( \rho(X) > 0 \)
        and \( \rho \in \rejectidb{\{X\}} \).
        Conversely, let \( \rho \in \rejectidb{\{X\}} \), then \( \rho \neq \rho_{\bot} \) and if \( \rho = \rho_{\top} \)
        then it is already contained in all \( \rejectidb{U} \), so we can focus on \( \rho \in \mathcal{R} \). 
        By Definition \( c := \rho(X) > 0 \), since \( \rho \in \mathcal{R} \) is Lipschitz continuous, there exists
        \( \delta > 0 \), s.t. \( \norm{Y - X}_{\infty} < \delta \implies |\rho(Y) - \rho(X)| < \frac{c}{2} \),
        note that if \( \rho(Y) \leq 0 \) then from \( c = \rho(X) > 0 \), the absolute value can not be smaller than \( \frac{c}{2} \)
        (in fact no smaller than \( c \)), so \( \rho(Y) > 0 \) for all \( Y \in B_{\delta}(X)\).
        But this means that \( \rho \) is strictly positive on some open neighborhood of \( X \), hence \( \rho \in \cup_{x \in U} \rejectidb{U} \).
    \end{proof}
    Let \( \textbf{BDLat} \) denote the category of bounded distributive lattices.
    \begin{proposition} \label{sheafification}
        %: \text{im } \presh \to \Pi \presh_x
        There is a monomorphism \( \text{sh} : \presh \to \mathcal{F} \) of presheaves, from the presheaf \( \presh \) of Theorem \ref{rmpresheaf}
        to a sheaf \( \mathcal{F} : \textbf{Open}(\mathcal{X}) \to \textbf{BDLat} \) on \( \mathcal{X} \), given by
        \[
            \mathcal{F}(U) \mapsto \left \{ (\sigma_X)_{X \in U} \in \Pi_{X \in U} \presh_X | (\sigma_X)_{X \in U} \text{ satisfies } (*)  \right \},
        \]
        where \( (*) \) is the property that for every \( X \in U \) there exists some open neighborhood \(V \subseteq U\) of \( X \) and a section \( [\rho] \in \presh(V) \)
        s.t. for all \( Y \in V \), it holds \( \sigma_Y = \varphi_{\{Y\},V}([\rho]) \).
    \end{proposition}
    \begin{proof}
        %Use Lemma 6.17.5 of \cite{stackprojectsheafification} and 
        The categorical product of distributive bounded lattices coincides with the Cartesian product \cite{catprodlattice}
        of them, so \( \prod_{X \in U} \presh_X \) is the lattice where \( \wedge, \vee \) are defined componentwise on tuples
        \( ([\rho]_{\presh_X})_{X \in U} \). By the \( (\ast) \) property of \( (\sigma_X)_{X \in U}, (\delta_X)_{X \in U} \in \mathcal{F}(U) \) 
        we have for any \( X \in U \) open neighborhoods \( V_{\sigma}, V_{\delta} \) over which for each \( Y \in V_{\sigma}, Z \in V_{\delta} \) it holds 
        \begin{align*}
            \sigma_Y = \varphi_{\{Y\},V_{\sigma}}([\rho]) = [\rho]_{\rejectidb{Y}} \\
            \delta_Z = \varphi_{\{Z\},V_{\delta}}([\eta]) = [\eta]_{\rejectidb{Z}}
        \end{align*}
        Since this holds over the whole neighborhood we bend the notation a little and write \( \sigma|_{V_{\sigma}} = [\rho]_{\rejectidb{V_{\sigma}}} \),
        \( \delta|_{V_{\delta}} = [\eta]_{\rejectidb{V_{\delta}}} \).
        By construction \( X \in V_{\sigma} \cap V_{\delta} \), so we have an open non-empty intersection \( V' \) on which we can project 
        \( [\rho] \), \( [\eta]_{} \) and take the meet and join there. In the succeeding we express this with the meet but the argument with the join is identical.
        Therefore locally one can express the meet as:
        \begin{align} \label{eqwedgeinsheaf}
            (\sigma_X)_{X \in V'} \wedge (\delta_X)_{X \in V'} = 
            \varphi_{V', V_{\sigma}}([\rho]) \wedge \varphi_{V', V_{\sigma}}([\eta]) = [\eta \wedge \rho]_{V'}.
        \end{align}
        Let \( (\gamma_{X})_{X \in U} := (\sigma_X)_{X \in U} \wedge (\delta_X)_{X \in U}\), to verify that \( \mathcal{F}(U) \)
        is a sublattice of \( \prod_{X \in U} \presh_X \), 
        we need to show that \( (\gamma_{X})_{X \in U} \) still satisfies \( (\ast) \), but this follows from the previous, since for each \( X \in U \) 
        we constructed an open \( V' \) on which \( (\gamma_{X})_{X \in U} \) restricts for each \( Y \in V' \) to the expression from Equation \ref{eqwedgeinsheaf}.
        Any sublattice of a distributive lattice is also distributive and the bounds are of course
        \(\top := ([\top]_{\presh_X})_{X \in U}, \bot := ([\bot]_{\presh_X})_{X \in U} \).
        \newline
        \newline
        We will proceed with verifying the monomorphism but leave the sheaf property of \( \mathcal{F}(U) \) to the reader (who alternatively can convince
        themself of it by looking up the official term of this subject, provided after the proof).
        To be a morphism the following diagram should commute for \( U \subseteq \mathcal{X} \) and open \( V \subseteq U \):
        \[
        \begin{tikzcd}[row sep=large, column sep=large]
            \presh(U) \arrow[r, "\text{sh}"] \arrow[d, "\varphi_{V,U}"]& \mathcal{F}(U) \arrow[d, "\psi_{V,U}"] \\
            \presh(V) \arrow[r, "\text{sh}"] & \mathcal{F}(V),
        \end{tikzcd}
        \]
        where \( \psi_{V,U} \) is the restriction in \( \mathcal{F}(U) \), given by \( ([\rho])_{\presh_X})_{X \in U} \mapsto ([\rho])_{\presh_X})_{X \in V} \).
        %dropping all indices \( X \in U \setminus V \) of \( (\sigma_X)_{X \in U} \).
        Take \( [\rho] \in \presh(U) \) , s.t. we restrict
        \( [\rho] \) onto \( \presh(V) \), via \( \varphi_{V,U} : \extrm / \rejectidb{U} \to \extrm / \rejectidb{V} \).
        Apply then the map to get \( \text{sh} \circ \varphi_{V,U}([\rho])) = ([\rho]_{\presh_X})_{X \in V} \),
        but pointwise this is simply \( \varphi_{\{X\},V} : \extrm / \rejectidb{V} \to \extrm / \rejectidb{\{ X \} } \),
        so we have the equality 
        \[
            \varphi_{\{X\},V} \circ \varphi_{V,U}([\rho]) = (\text{sh} \circ \varphi_{V,U}([\rho])))_X,
        \]
        where the subscript \( X \) denotes the component of \( \text{sh} \circ \varphi_{V,U}([\rho]) \in \mathcal{F}(V) \) at \( X \).
        Since \( \presh \) is a presheaf the LHS simplifies:
        \[
         \varphi_{\{X\}, U}([\rho]) = (\text{sh} \circ \varphi_{V,U}([\rho])))_X
        \]
        which is precisely \( (\psi_{V,U} \circ \text{sh}([\rho]))_X \) and as \( X \) was arbitrary, the diagram commutes.
        That the morphism is injective follows, because if
        we have \( ([\rho_1]_{\presh_X})_{X \in U} = ([\rho_2]_{\presh_X})_{X \in U} \), we get for each \( X \) an \( \eta_X \in \rejectidb{\{X\}} \)
        s.t. \( \rho_1 \wedge \eta_X = \rho_2 \wedge \eta_X \) and similar to the proof from Theorem \ref{rmpresheaf} for separatedness we get 
        an \( \eta = \bigvee_{X \in U} \eta_X \) in \( \rejectidb{U} \) which satisfies \( \rho_1 \wedge \eta = \rho_2 \wedge \eta \).
    \end{proof}
    The previous map is in fact the sheafification of \( \presh \) whose concatentation with the map is denoted by \( \mathcal{F} \), which satisfies the universal property that for any sheaf \( \mathcal{G} \) and map of presheaves \( \psi : \presh \to \mathcal{G} \), there is a unique map \( \psi^\sharp : \mathcal{F} \to \mathcal{G} \) s.t. 
    \( \psi^\sharp \circ \text{sh} = \psi \).
    Without going full abstract nonsense (if we haven't done so already),
    it resolves the obstruction of gluing our local sections to a global section. This introduces objects which we will refer to as local risk measures.
    \begin{definition} \label{deflocalrm}
        A local risk measure \( \rho : U \to \extreals \) with open \( U \subseteq \mathcal{X} \) is a map 
        s.t. for each \( X \in U \), there is some open \( V \subseteq U \) on which \( \rho \)  restricts to an element in \( \extrm \).
    \end{definition}
    \begin{remark} \label{analoglocalrm}
        Definition \ref{deflocalrm} already looks similar to the sections of the sheaf from Proposition \ref{sheafification}
        in regard to local behavior. That they in fact specify the same object is however not immediate.
        First one needs to define the local analogues \( \extrmloc, \rejidbloc{U}, \accidbloc{U} \) verify 
        that \( \extrmloc \) is indeed a distributive bounded lattice which is in addition complete, JID/MID
        and that \( \rejidbloc{U} \) is a filter. This task is left to the reader, if they wish to do so.
        The actual isomorphism is something we will cover later.
    \end{remark}
    %(TODO it is not immediate how such objects can be used to match sections in the sheaf from \ref{sheafification}) \newline
    As concrete example why we need such objects for a sheaf and a counterexample that \( \presh \) from Theorem \ref{rmpresheaf}
    is only a presheaf, take any disjoint open sets \( U_1, U_2 \) in \( \mathcal{X} \) s.t. \( \rho_1 = \text{VaR}_{\lambda} > 0 \) on \( U_1 \) and 
    \( \rho_2 = \mathbb{E}[-X] > 0 \) on \( U_2 \), so \( \rho_1 \in \rejectidb{U_1}, \rho_2 \in \rejectidb{U_2} \), since \( U_1 \cap U_2 = \emptyset \) and \( \rho_1, \rho_2 \) 
    restrict to \( \rho_{\top} \) on \( \rejectidb{\emptyset} \), there must be a risk measure \( \rho \in \rejectidb{U_1 \cup U_2} \), if \( \presh \) would be a sheaf, 
    restricting to \( \rho_i \) on \( U_i \), but this seems impossible for non trivial \( \Omega, \mathbb{P} \). There are also more realistic examples:
    \begin{example}
        Let \( \rho_1 = \text{ess sup} - X \), the worst case risk measure. We will define some intermediate 
        risk measures for the actual counter example:
        \begin{align*}
            \eta_i(X) = \begin{cases}
                \text{ess sup} (- X) + i\cdot 100, \text{ for } i \leq 2 \\
                 \expv{-X} + i \cdot 1000, \text{ else}
            \end{cases}
        \end{align*}
        we combine them to risk measures 
        \begin{align*}
            \rho_2 &= \rho_1 \wedge (\eta_1 \vee \eta_5) \\
            \rho_3 &= (\rho_1 \vee \eta_3) \wedge (\eta_2 \vee \eta_4)
        \end{align*}
        with corresponding acceptance sets
        \begin{align*}
            \mathcal{A}_{\eta_i} = \begin{cases}
                \{ X \in \mathcal{X} | X \geq i \cdot 100 \text{ a.s. } \}, \text{ for } i \leq 2 \\
                \{ X \in \mathcal{X} | \expv{X} \geq i \cdot 1000 \}, \text{ else}
            \end{cases}
        \end{align*}
        and \( \mathcal{A}_{\rho_2}, \mathcal{A}_{\rho_3} \) the corresponding intersections and unions of them (recall \ref{acceptanceoperations}):
        \begin{align*}
            \mathcal{A}_{\rho_2} &= \mathcal{A}_{\rho_1} \cup (\mathcal{A}_{\eta_1} \cap \mathcal{A}_{\eta_5}) \\
            \mathcal{A}_{\rho_3} &= \left(\mathcal{A}_{\rho_1} \cap \mathcal{A}_{\eta_3} \right) \cup \left(\mathcal{A}_{\eta_2} \cap \mathcal{A}_{\eta_4} \right)
        \end{align*}

        Let \( U := B_{100}(1) \) then all \( X \in U \) are naturally in \( \mathcal{A}_{\rho_1} \) 
        and then also in \( \mathcal{A}_{\rho_2} \). By assumption \( 99 < \expv{X} < 101  \)
        so \( X \notin \mathcal{A}_{\eta_5} \), therefore \( \rho_2 \) coincides with \( \rho_1 \) on \( U \).
        Further \( X \notin \mathcal{A}_{\eta_3} \) and \( X \notin \mathcal{A}_{\eta_4} \), and therefore \( X \notin \mathcal{A}_{\rho_3} \)
        implying \( \rho_3 \in \rejectidb{U} \).
        Observe, that \( \rho_1 \wedge \eta_1 = \rho_1 \), and therefore \( \rho_1 \wedge (\eta_1 \vee \eta_5) = \rho_1 \).
        This enables the following set of equivalences:
        \begin{align*}
            \rho_1 \wedge (\eta_2 \vee \eta_4) &= \rho_1 \wedge (\eta_2 \vee \eta_4) \\
            \rho_1 \wedge (\rho_1 \vee \eta_3) \wedge (\eta_2 \vee \eta_4) &= \rho_1 \wedge (\eta_1 \vee \eta_5) \wedge (\eta_2 \vee \eta_4) \\
            \rho_1 \wedge \rho_3 &= \rho_1 \wedge (\rho_1 \vee \eta_3) \wedge (\eta_1 \vee \eta_5) \wedge (\eta_2 \vee \eta_4) \\
            \rho_1 \wedge \rho_3 &= \rho_2 \wedge \rho_3
        \end{align*}
        In particular, this means that \( [\rho_2] = [\rho_1] \) in \( \presh(U) \), but they can't be extended to a global risk measure
        in \( \extrm = \presh(\mathcal{X}) \), since they are distinct.
        The intuition behind the individual risk measures is that the viewpoint \( \rho_2 \) is also content with risky positions which have a high return 
        as opposed to \( \rho_1 \) which only takes on safe positions. \( \rho_3 \) requires a higher return then \( \rho_1 \) on safe positions, but is more lenient
        than \( \rho_2 \) for risky ones.
    \end{example}

    With the Definitions from Remark \ref{analoglocalrm} we now show the direct approach to our sheaf.
    \begin{theorem}
        There is an isomorphism \( \Psi : \extrmloc / \rejidbloc{U} \to \mathcal{F}(U) \) of bounded distributive lattices
        for each \( U \subseteq \mathcal{X} \).
    \end{theorem}
    \begin{proof}
        As an intuition of what is to come, one might refresh their repertoire on partitions of unity.
        %We will utilize a common trick in the construction of sheafs, which is conceptually the same as a partition of unity.
        Let \( \psi \in \extrmloc \) be a local risk measure and define \( \Psi([\psi]) \) as the section \( (\sigma_X)_{X \in U} \in \mathcal{F}(U) \)
        s.t. for each \( X \in U \) and corresponding neighborhood \( V \subseteq \) of \( X \) on which \( \psi \) restricts to a standard risk measure \( \rho \in \extrm \)
        we have \( \sigma_Y = \varphi_{\{Y\}, V}([\rho]) \). This already fullfills the \( (\ast) \) condition of Proposition \ref{sheafification}.
        \newline
        \newline
        We will show well-definedness of \( \Psi \) with injectivity at once.
        Take \( \psi_1, \psi_2 \in \extrmloc \) and assume \( \Psi([\psi_1]) = \Psi([\psi_2]) \) in \( \mathcal{F}(U) \),
        implying for each \( Z \in U \) and corresponding neighborhood \( V \) and restrictions \( \rho_1, \rho_2 \)
        of \( \psi_1, \psi_2 \) some \( \gamma_Z \in \rejectidb{V_Z} \) s.t. \( \gamma_Z \wedge \psi_1 = \gamma_Z \wedge \psi_2 \)
        and we can construct an open cover \( \{V_Z\}_{Z \in U} \) of \( U \) with them.
        As \( \mathcal{X} = L^\infty(\Omega, \mathcal{F}, \mathbb{P}) \) is a metric space, so is every open subset of it, 
        whence every \( U \) is paracompact yielding some locally finite refinement \( \{W_Z\}_{Z \in U'} \) of \( \{V_Z\}_{Z \in U} \).
        Define 
        \begin{align} \label{patchedlocalrm}
            \gamma(X) := \inf_{\{Z | X \in W_Z \}} \gamma_Z(X).
        \end{align}
        since \( X \) is covered by at least one \( W_Z \) the set over which we take infimum is non-empty 
        and for every included \( W_Z \) it holds \( \gamma_Z(X) > 0 \). Since \( \{W_Z\}_{Z \in U'} \) is locally finite,
        only finitely many \( W_Z \) are taken into account and therefore this is even a minimum, so \( \gamma(X) > 0 \).
        Since the set of \( W_Z \) over which we take the infimum is constant over some neighborhood of \( X \), by Definition
        of local finiteness, \( \gamma \) is a risk measure there and by arbitrariness of \( X \), it follows \( \gamma \in \rejidbloc{U} \).
        \iffalse
        Since for every \( X \in U \) there is some open neighborhood \( V \), which intersects only finitely many \( W_Z \), call them active, 
        which is again an open, we can focus on the intersection of \( V \) and all active \( W_Z \), which by construction is non-empty. 
        Here \( \gamma \) restricts just to a regular risk measure, namely \( \bigwedge_{Z^\ast} \gamma_{Z^\ast}(X) \), where \( Z^\ast \)
        correspond to the active \( W_Z \). As we could do this for each \( X \in U \), we have .
        \fi
        Moreover, we have 
        \[
            \psi_1 \wedge \gamma = \psi_1 \wedge \left (\bigwedge_{Z \in U'} \gamma_{Z} \right) = \bigwedge_{Z \in U'} \psi_1 \wedge \gamma_{Z}
        \]
        and in any specific \( W_Z \) it holds \( \psi_1 \wedge \gamma_Z = \psi_2 \wedge \gamma_Z \). Therefore the term rewrites to
        \[
            \bigwedge_{Z \in U'} \psi_2 \wedge \gamma_{Z} = \psi_2 \wedge \gamma,
        \]
        yielding \( [\psi_1] = [\psi_2] \) in \( \extrmloc / \rejidbloc{U} \).
        \newline
        \newline
        Let \( (\sigma_X)_{X \in U} \in \mathcal{F}(U) \) and extract a locally finite cover \( \{W_Z\}_{Z \in U'}\) from the cover induced by the defining \( V \)
        of the \( (\ast) \) property. 
        %On each \( W_Z \) we get at least one representative \( \rho \in \extrm \) with \( \sigma_Y = \varphi_{\{Y\}, V}([\rho]) \)
        %for \( Y \in W_Z \).
        Use the axiom of choice to pick a representative from each class \( [\rho_Z] \in \mathcal{F}(W_Z) \) with \( [\rho_Z] = \sigma |_{W_Z} \) and define
        \[
            \psi(X) := \inf_{\{Z | X \in W_Z\}} \rho_Z(X),
        \]
        deduce analogously to the construction of Equation \ref{patchedlocalrm}, that \( \psi \in \extrmloc / \rejidbloc{U} \).
        Suppose \( \sigma|_{V_X} = [\rho] \) with \( \rho \in \rejectidb{V_X} \), we then have for each \( Y \in V_X \) either \( \rho(Y) \leq 0 \)
        or \( \rho(Y) > 0 \). In the first case, also \( \rho_Z(X) = \rho(Y) \) for all \( \rho_Z \) with \( W_Z \cap V_X \neq \emptyset \), 
        since they are in the same equivalence class, therefore \( \psi(Y) = \rho(Y) \) and \( [\psi]_{Y} = \sigma_Y \). Otherwise all \( \rho_Z(X) > 0 \),
        but not necessarily with the same value. Regardless, it follows \( [\psi]_{Y} = [\top] = \sigma_Y \), and therefore \( [\psi]|_{V_X} = \sigma \mid_{V_X} \).
        By arbitrariness of \( X \), we conclude with \( \Psi(\psi) = \sigma \).
    \end{proof}

    \begin{corollary}
        There exists a sheaf \( \mathcal{F} : \textbf{Open}(\mathcal{X}) \to \textbf{BDLat} \) on \( \mathcal{X} \),
        which we call the rejection risk sheaf of \( \mathcal{X} \).
    \end{corollary}
    \begin{proof}
        This is just a restatement of Proposition \ref{sheafification}.
    \end{proof}
    %As noted before the presheaf \( \presh \) and \( \mathcal{F} \) produce for general \( U \subseteq \mathcal{X} \) not the same 
    Local risk measures are interesting in their own right, since they could reflect the fact that locally the behavior of an actor
    appear rational or consistent, but globally they contradict themselves. Meaning they are willing to take on unnecessary risk in some situations.
    But there's a special case, where we can guarantee a clear risk profile globally.
    %There's a special case
    %with the 
    %There is a special case when local risk measures are forced to be a global risk measure.
    \begin{proposition}
        Let \( U \subseteq \mathcal{X} \) be convex, then any local risk measure \( \psi : U \to \extreals \) coincides with a global risk measure
        on \( U \).
    \end{proposition}
    \begin{proof}
        Define for each \( X \in U \), the function \( f(m) = \psi(X + m) + m \) and let \( V \) be the corresponding neighborhood of \( X \)
        in which \( \phi \) restricts to some member of \( \extrm \). In \( V \) we have \( \partial_m f(m) = 0 \) therefore \( \psi \) is locally constant.
        As any convex \( U \) is connected, a locally constant function is globally constant, hence \( \psi(X + m) = \psi(X) - m \) for all \( m \in \mathbb{R} \)
        s.t. \( X + m \in U \).
        \newline
        \newline
        Let \( X, Y \in U \) s.t. \( X \leq Y \) and by definition the path \( Z_\lambda := \lambda Y + (1 - \lambda)X \) is contained in \( U \) for all \( \lambda \in [0, 1] \).
        For fixed \( \omega \in \Omega \) we have \( \partial_\lambda Z_\lambda(\omega) = Y(\omega) - X(\omega) \geq 0 \), so \( Z_\lambda \) is a.s. inreasing.
        Let \( \{ V_\lambda \}_{\lambda \in [0, 1]} \) be an open covering of the path \( Z_\lambda \) in its subspace topology with the \( V_\lambda \) 
        corresponding to the restriction neighborhoods of \( \psi \). Since the map \( \lambda \mapsto Z_\lambda \) is a.s. continuous, we obtain some open cover
        \( \{I_\lambda\}_{\lambda \in [0, 1]} \) of \( [0, 1] \) via its inverse. By compactness of \( [0, 1] \) this has some finite subcover \( \{I_\lambda\}_{\lambda \in \Lambda} \)
        and in each interval \( Z|_{I_\lambda} \) it holds for \( \lambda_1 \leq \lambda_2 \), \( \varphi(Z_{\lambda_1}) \geq \varphi(Z_{\lambda_2}) \). Since 
        the restrictions of \( \varphi \) coincide on overlaps of the \( Z_{I_\lambda} \), we obtain \( \varphi(Z_0) \geq \varphi(Z_1) \), so \( \varphi \) is a global risk measure
        on \( U \).
    \end{proof}

    In the following sections, we will discuss the meaning of robust representations from a lattice point of view and 
    by doing so stay very close to the underlying geometric interpretation. However, our sheaf construction won't re-enter the stage anymore. 
    The reader might wonder what they takeaway of this should be then. First, there are not many existing sheafs in literature 
    which map into the category of lattices, so this is a new addition to this collection. Secondly, it shows that modern tools of algebraic geometry
    can be translated to non-polynomial settings aside from the usual algebraic number theory. The obstruction to study the geometry of acceptance and rejection sets, 
    here is essentially that the ideals and filters are a very one-sided perspective. In algebraic geometry we have the luxury to be able to still evaluate sections on the underlying space,
    while for our quotient lattices this was only possible for the elements, which were accepted. Essentially one would need to consider two sheafs, 
    one for the rejection filters and another for acceptance ideals (which have not built here, but the motivated reader can translate the previous proofs to this setting) 
    to get the full picture, but the author perceived this as adding more unnecessary complexity than insight. A better behaved version might be the sheaf of local risk measures
    that is \( U \mapsto \extrmloc(U) \), where \( \extrmloc(U) = \{ \psi|_U | \psi \in \extrmloc \} \), but this is out of the scope of this thesis.
\end{subsection}
